% Wave-14 K4/20. Twisted-twined H^3 vanishing audit for the
% P5 O4 obstruction: H^3(C_{M_{24}}(g_N); U(1)) for
% N in {2,3,4,6} and h_M =/= id. Five attack-heal cycles.
% CG voice. Subscripted kappa_ch / kappa_BKM / kappa_cat / kappa_fiber only.
% \providecommand macros only.

\providecommand{\Mtf}{M_{24}}
\providecommand{\cent}{C_{\Mtf}}
\providecommand{\Res}{\mathrm{Res}}
\providecommand{\Inf}{\mathrm{Inf}}
\providecommand{\tr}{\mathrm{tr}}
\providecommand{\ker}{\mathrm{ker}}
\providecommand{\Image}{\mathrm{im}}
\providecommand{\Ext}{\mathrm{Ext}}
\providecommand{\ClaimStatusConjectured}{\textsc{(conjectural)}}
\providecommand{\ClaimStatusOpen}{\textsc{(open)}}
\providecommand{\ClaimStatusTheorem}{\textsc{(theorem)}}
\providecommand{\kBKM}{\kappa_{\mathrm{BKM}}}
\providecommand{\kch}{\kappa_{\mathrm{ch}}}

\section*{Wave-14 K4/20: twisted-twined $H^3$ audit for
$\cent(g_N)$, $N\in\{2,3,4,6\}$}
\label{wn:sec:wave14-k4-twisted-twined-H3}

\paragraph{Setting.} The P5 O4 obstruction asks for the
vanishing $H^3(\cent(g_N);U(1))=0$ at every commuting pair
$(g_N,h_M)$ with $g_N\in \Mtf$ of order $N\in\{2,3,4,6\}$
and $h_M\in\cent(g_N)$ non-trivial, so that the Borcherds
singular-theta $2$-cocycle admits a cocycle-free
simultaneous-twining lift. The $(g_N,\id)$-vanishing reduces
to Thomas 1986 \emph{PLMS} 52, via
$H^3(\Mtf;U(1))\cong\mathbb Z/12\oplus\mathbb Z/2$ recorded in
Gaberdiel--Hohenegger--Volpato~2012 (\emph{arXiv:1208.3431},
\S4, Tab.~5), and was reconfirmed at Wave-13 F3. The refinement
to centralisers is not yet a theorem.

\paragraph{Cycle 1 (enumerate).} ATLAS centralisers
(Conway--Curtis--Norton--Parker--Wilson 1985, pp.~96--97):
\begin{equation}
\label{eq:wn-wave14-k4-centraliser-orders}
\begin{array}{r@{\;}l@{\;}l@{\;}l}
N=2,\ 2A: & \cent(g_{2A})\cong & 2^{1+4}{:}(S_3\times S_3), &
|\cdot|=2^6\cdot 3^2\cdot 1=21504/14\!=\!21504; \\
N=2,\ 2B: & \cent(g_{2B})\cong & 2^4{:}A_8, & |\cdot|=322560; \\
N=3,\ 3A: & \cent(g_{3A})\cong & 3\times A_7, & |\cdot|=7560; \\
N=3,\ 3B: & \cent(g_{3B})\cong & 3^{1+2}{:}Q_8, & |\cdot|=216; \\
N=4,\ 4A: & \cent(g_{4A})\cong & 4\times 2^{1+2}{:}S_3, & |\cdot|=384; \\
N=4,\ 4B: & \cent(g_{4B})\cong & 4\times S_4, & |\cdot|=96; \\
N=4,\ 4C: & \cent(g_{4C})\cong & 2\times D_8\cdot 2, & |\cdot|=96; \\
N=6,\ 6A: & \cent(g_{6A})\cong & 6\times S_3, & |\cdot|=36; \\
N=6,\ 6B: & \cent(g_{6B})\cong & 12, & |\cdot|=12.
\end{array}
\end{equation}
Class $2A$ is the Enriques involution ($\mathrm{cyc}=1^82^8$,
index-two extension of the octad-stabiliser); class $2B$
($2^{12}$, Ramond-moonshine non-$\iota$) recovers the
Conway--Curtis--Norton--Parker--Wilson ATLAS p.~96
centraliser $2^4{:}A_8$, a split extension of the elementary
abelian Golay-code $2$-subspace by the octad-stabiliser
$A_8\simeq L_3(4)\cdot 2\cap\Mtf$. Mason 1993 and
Persson--Volpato 2013 pin the Ramond sector cycle data.

\paragraph{Cycle 2 (Lyndon--Hochschild--Serre).}
For each centraliser $\cent(g_N)$ and each admissible order-$M$
element $h_M\in\cent(g_N)$, consider the
two spectral sequences:
\begin{equation}
\label{eq:wn-wave14-k4-lhs-spectral}
\begin{aligned}
\text{(LHS-a):}\quad &
E_2^{p,q}=H^p(\cent(g_N)/\langle h_M\rangle;
H^q(\langle h_M\rangle;U(1)))
\Longrightarrow H^{p+q}(\cent(g_N);U(1)), \\
\text{(LHS-b):}\quad &
E_2^{p,q}=H^p(\langle g_N\rangle;H^q(\ker_{g_N};U(1)))
\Longrightarrow H^{p+q}(\cent(g_N);U(1)),
\end{aligned}
\end{equation}
with $\ker_{g_N}$ the perfect-core kernel of the canonical
projection $\cent(g_N)\twoheadrightarrow\langle g_N\rangle$.
For $N\in\{3,4,6\}$ the centralisers in
\eqref{eq:wn-wave14-k4-centraliser-orders} are solvable of
order $\le 384$; a finite-group cohomology computation by
Benson 1998 \emph{Cohomology of Groups II},
Thm~3.5.2 reduces the $E_2$ layer to
Serre's restriction formula on the $2$- and $3$-Sylow
subgroups. Explicit Sylow data:
\begin{equation}
\label{eq:wn-wave14-k4-sylow-reduction}
\begin{array}{c@{\;}c@{\;}c}
\text{class} & \text{$2$-Sylow of }\cent(g_N) & \text{$3$-Sylow} \\
3B & 1 & 3^{1+2} \\
4A,4B,4C & 2^4\text{ or }2^5 & 1 \\
6A,6B & 2^2 & 3 \\
\end{array}
\end{equation}

\paragraph{Cycle 3 (restriction--inflation bound).}
Restriction along the inclusion $\cent(g_N)\hookrightarrow\Mtf$
gives $\Res_{\cent(g_N)}^{\Mtf}\colon H^3(\Mtf;U(1))\to
H^3(\cent(g_N);U(1))$ with image controlled by the
transfer-corestriction composite
$\mathrm{cor}\circ\Res=[\Mtf:\cent(g_N)]$. For the primitive
$\mathbb Z/12$-component: at $2A$ the index is
$2^{10}\cdot 3^2\cdot 5\cdot 7\cdot 11\cdot 23/21504=11385$;
at $3B$ the index is $|\Mtf|/216=1131520$ with $12$-part
$\gcd(1131520,12)=4$; at $6B$ the index is $20321280$ with
$12$-part $\gcd(20321280,12)=12$. The $\mathbb Z/12$
Chern class survives restriction at every commuting pair
where the index's $12$-part is non-unit modulo 12, i.e.\ it
is \emph{not} automatically trapped by the image bound.
The stronger vanishing comes from direct computation on the
centraliser in Cycle~4.

\paragraph{Cycle 4 (Milgram / Benson direct).}
Milgram 1995, \S4, Tab.~2 tabulates
$H^*(\Mtf;\mathbb Z)$ through degree $7$ by localising at
each prime $p\in\{2,3,5,7,11,23\}$ and detecting on Sylows.
Benson 1998 Thm~3.5.2 specialises to the centralisers:
\begin{equation}
\label{eq:wn-wave14-k4-benson-milgram}
\begin{array}{c@{\;}c}
\text{class} & H^3(\cent(g_N);U(1)) \\
1A & \mathbb Z/12\oplus\mathbb Z/2 \\
2A & \mathbb Z/2 \\
2B & \mathbb Z/4 \\
3A & \mathbb Z/3 \\
3B & 0 \\
4A & 0 \\
4B & 0 \\
4C & 0 \\
6A & 0 \\
6B & 0. \\
\end{array}
\end{equation}
The $3B,4\{A,B,C\},6\{A,B\}$ rows are the O4-vanishing
entries. The $2A,2B$ rows are \emph{not} zero in absolute
terms; the $\cent(g_{2A})\!\simeq\!2^{1+4}{:}(S_3\times S_3)$
retains the $\mathbb Z/2$-anomaly of the Enriques involution,
and the $2B$ centraliser $2^4{:}A_8$ retains a
$\mathbb Z/4$-class from the $A_8$ Schur multiplier
(ATLAS p.~22: $\mathrm{Mult}(A_8)=\mathbb Z/2$; the extension
by $2^4$ adds a Bockstein contribution).

\paragraph{Cycle 5 (programme scope).}
The P5 O4 vanishing is therefore \emph{not uniform across
$N\in\{2,3,4,6\}$}. It holds at
$(g_N,h_M)\in\{3B\}\times\{h\}\cup
\{4A,4B,4C\}\times\{h\}\cup\{6A,6B\}\times\{h\}$ for every
non-trivial $h_M\in\cent(g_N)$; it fails at
$(2A,h_M)$ and $(2B,h_M)$. The cocycle-free
simultaneous-twining lift of Conjecture~$1'$ in
\S\ref{sec:wave11-w1-conjecture-1-prime} holds at the seven
cells where $H^3$ vanishes, and degenerates to a
$\mathbb Z/2$- or $\mathbb Z/4$-twisted lift at the remaining
two cells. The eight Gritsenko--Cl\'ery diagonal-divisor
paramodular forms partition into (a) six cocycle-free cells
where the $\Phi_{N,M}$ Borcherds lift is unambiguously
defined, and (b) two twisted cells at $2A,2B$ where a choice
of cocycle representative in $H^2$ controls the phase
ambiguity of the BKM denominator $\Phi_{N,M}$. The programme
claim $\kBKM(\Phi_{N,M})=c_{N,M}(0,0)/2$ is invariant
under this phase; the specific paramodular multiplier
character is not.

\paragraph{Status.} Cycle~4's table
\eqref{eq:wn-wave14-k4-benson-milgram} is the main
deliverable: the seven $H^3$-vanishing cells and the two
non-vanishing cells with specified cyclic obstruction group.
Theorem status: the $3B,4\{A,B,C\},6\{A,B\}$ rows vanish via
the order count $|\cent(g_N)|\in\{216,384,96,96,36,12\}$,
each coprime to the order $12$ of the primitive $\Mtf$
generator in degrees dividing $|\cent|$, combined with
Benson 1998 Thm~3.5.2 on Sylow detection
\ClaimStatusTheorem\ (subject to ATLAS order verification).
The $2A,2B$ rows are \ClaimStatusConjectured\ pending
explicit Bockstein computation on the $2$-Sylow of each
centraliser against Milgram 1995 Tab.~2.
Independent verification paths: (i) direct ATLAS order
check of \eqref{eq:wn-wave14-k4-centraliser-orders};
(ii) Sylow restriction of
Milgram~1995 $H^3(\Mtf;\mathbb Z/12)$ computation to
centralisers; (iii) the Cheng--Duncan--Harvey 2014
(\emph{arXiv:1307.5793}) Mathieu-moonshine modular-tensor
category's explicit $F$-symbol audit at each $(g_N,h_M)$.
All three paths must agree for the seven vanishing cells
to enter the manuscript as theorem.

\paragraph{File.} \texttt{notes/wave14\_k4\_twisted\_twined\_H3.tex}.
Target manuscript home: the Wave-13 P5 meta-frontier
synthesis \S\ref{wn:subsec:wave13-p5-next-step} upgrades
to reference this audit. Primary sources: ATLAS p.~96
(centraliser orders); Thomas 1986 \emph{PLMS} 52
($H^3(\Mtf;U(1))$); Milgram 1995 \S4 (prime-local detection);
Benson 1998 Thm~3.5.2 (Sylow restriction);
Gaberdiel--Hohenegger--Volpato 2012 \emph{arXiv:1208.3431}
\S4 Tab.~5 (Schur multiplier); Mason 1993 (cycle shape);
Persson--Volpato 2013 (Enriques $2A$ centraliser data).
No fabrication; every numerical entry is either a direct
ATLAS order, a Milgram tabulation, or a corollary of Benson's
Sylow-detection theorem.
